\documentclass[10pt,a4paper]{amsart}
\usepackage[margin=19mm]{geometry}
\usepackage{amsmath,amssymb,mathtools}
\usepackage[scr=rsfs,cal=euler]{mathalfa}
\usepackage{xfrac}
\usepackage[unicode,hidelinks,bookmarks=false]{hyperref}
\newcommand{\ie}{i.e.\ }
\newcommand{\cf}{cf.\ }
\newcommand{\resp}{respectively\ }
\newcommand{\Subsection}{\subsection}
\providecommand{\nobreakdash}{\nobreak\hbox{-}}
\setlength{\emergencystretch}{2em}
\multlinegap0pt
\usepackage{xcolor}
\usepackage[shortlabels]{enumitem}
\usepackage{amssymb}
\usepackage{mathtools}
\newtheorem{prop}{Proposition}
\newcommand{\qop}{\triangleleft}
\title{A condition of admissibility for generalized Alexander quandles}
\author{Katsunori Arai, Keisuke Himeno, Ryoya Kai, Yuko Ozawa}
\date{Source: arXiv:2608.25345v1 (CC BY 4.0)}
\begin{document}
\maketitle

\noindent\emph{Source and licence.} A condition of admissibility for generalized Alexander quandles. Version arXiv:2608.25345v1 (CC BY 4.0), distributed by arXiv under the Creative Commons Attribution 4.0 International licence, http://creativecommons.org/licenses/by/4.0/, as recorded in the arXiv licence metadata for this exact version and shown on the abstract page https://arxiv.org/abs/2608.25345v1. Full licence notice: \texttt{licenses/arxiv-2608.25345-CC-BY-4.0.txt}.\par\smallskip

\noindent\textit{Editorial scope.} Counted unit: the article's prop prop:max\_anti\_GAlex (source0.tex lines 393--405). The article's complete original proof of it is delivered with it (source0.tex lines 406--439, one \texttt{proof} environment of 34 lines). No mathematics is written, repaired or re-derived: the statement and the proof are the article's own lines, in the article's own notation, and no retained line is substituted or re-flowed.  No further source-local context is needed: every cross-reference of every retained line resolves inside the two delivered spans.  Nothing was omitted from any retained span, so the package carries no omission record.  No retained line contains a citation, so the wrapper carries no bibliography and no bibliography entry.  No retained line contains a figure environment, an image-inclusion command or drawing code, so no image is delivered and the package declares no asset dependency.  Editorial formatting only, in addition: an AMS \texttt{amsart} preamble that restates the article's own declarations and macros used by the retained lines, namely the environments \texttt{prop} on separate counters, together with the standard packages those lines need.  The retained environments print in this document as Proposition 1 in order of appearance, whereas the article prints the same items with the numbering of its own full document.  The complete native source of the article as distributed in the arXiv e-print for this exact version is bundled unchanged as \texttt{sources/208569/source0.tex}.
\begin{prop}\label{prop:max_anti_GAlex}
  % A subset $T \subset X$ containing $g \in X$ is antipodal if and only if $T \subset Hg$.
  The following hold:
  \begin{enumerate}
    \item[$(1)$] For a subset $T \subset X$ containing $g \in X$,
    the set $T$ is antipodal if and only if $T \subset Hg$.

    \item[$(2)$] A maximal antipodal set of $X$ containing $g \in G$ is equal to the subset $Hg$.

    \item[$(3)$] The subset $Hg$ is a pole of $X$.
    % If two elements $x, y \in X$ satisfy $\eta_X(x) = \eta_X(y)$, then $xy^{-1} \in H$.
  \end{enumerate}
\end{prop}

\begin{proof}
  We first show $(1)$.
  We assume that a subset $T$ containing $g$ is antipodal.
  Let us take $x \in X$.
  Then, it satisfies $x = x \qop g = \sigma(x g^{-1}) g$,
  and hence $xg^{-1} \in H$.
  Thus, we have $x \in Hg$.
  This shows that $T \subset Hg$.
  To show the inverse, we prove that the subset $Hg$ is antipodal.
  For $hg, h'g \in H g$, we have
  \[
    hg \qop^{\pm} h'g = \sigma^{\pm 1}(hg(h'g)^{-1}) h'g = \sigma^{\pm 1}(h h'^{-1})h'g = h h'^{-1}h'g = hg,
  \]
  as desired.

  Next, we show $(2)$.
  To show this, it suffices to show that $Hg$ is a maximal antipodal set by $(1)$.
  Suppose that $T$ is a trivial subquandle of $X$ with $Hg \subset T$.
  Let us take $t \in T$.
  Since $T$ is a trivial quandle containing $g \in T$, it satisfies the following:
  \[
    t = t \qop g = \sigma(tg^{-1})g.
  \]
  Thus, we have $\sigma(tg^{-1}) = tg^{-1}$.
  This shows $tg^{-1} \in H$, and hence $t \in Hg$.
  Therefore, we conclude that $T = Hg$, as desired.

  Finally, we show $(3)$.
  For $hg, h'g \in Hg$, we have
  \begin{align*}
      (hg) \qop (h'g) = \sigma((hg)(h'g)^{-1})h'g = \sigma(hh'^{-1})h'g = hh'^{-1}h'g = hg,
  \end{align*}
  which completes the proof.
\end{proof}

\end{document}
